\title{Controlling Text-to-Image Diffusion by Orthogonal Finetuning}

\begin{abstract}
Large-scale text-to-image diffusion models demonstrate remarkable ability to produce photorealistic images from textual descriptions. An important open problem is how to effectively steer or control these powerful models to perform various downstream tasks. To address this, we propose a principled finetuning approach—Orthogonal Finetuning (OFT)—for adapting text-to-image diffusion models to specific downstream applications. Unlike prior methods, OFT can theoretically maintain the hyperspherical energy, which quantifies pairwise neuron relationships on the unit hypersphere. We find this characteristic is vital for preserving the semantic generation capabilities of text-to-image diffusion models. To enhance finetuning stability, we additionally introduce Constrained Orthogonal Finetuning (COFT), which enforces a radius constraint on the hypersphere. Specifically, we focus on two key finetuning tasks: subject-driven generation, aiming to create subject-specific images based on a few example images and a text prompt, and controllable generation, which enables the model to incorporate supplementary control signals. Empirical results demonstrate that our OFT framework surpasses existing techniques in terms of both generation quality and convergence speed.
\end{abstract}

\section{Introduction}

Recent advances in text-to-image diffusion models~\cite{saharia2022photorealistic,ramesh2022hierarchical,rombach2022high} have shown outstanding capabilities in generating high-fidelity images guided by textual input. Despite these strong outcomes, text prompts can still be vague and insufficient for providing detailed and precise control over the generated images. To overcome this limitation, we focus on two specific types of text-to-image generation tasks in this work:

\begin{itemize}[leftmargin=*,nosep]

    \item \textbf{Subject-driven generation}~\cite{ruiz2023dreambooth}: The objective is to produce images of a particular subject in varied settings, given only a handful of example images of that subject along with a textual prompt.
    \item \textbf{Controllable generation}~\cite{zhang2023adding,mou2023t2i}: The goal is to generate images conditioned on both a text prompt and an additional control signal (e.g., canny edges, segmentation maps).
\end{itemize}

Both tasks fundamentally reduce to the challenge of effectively finetuning text-to-image diffusion models without compromising their pretrained generative capabilities. We outline the key requirements for a successful finetuning approach as: (1) \emph{training efficiency}: minimizing the number of trainable parameters and training epochs, and (2) \emph{generalizability preservation}: maintaining high-fidelity and diverse generative outputs. Typically, finetuning is performed either by adjusting neuron weights with a small learning rate (\emph{e.g.}, \cite{ruiz2023dreambooth}) or by incorporating a minor component with re-parameterized neuron weights (\emph{e.g.}, \cite{hulora2022,zhang2023adding}). Although straightforward, neither approach guarantees the retention of pretrained generative performance. Additionally, there is a lack of theoretical insight into designing effective finetuning methods and determining appropriate hyperparameters such as the number of training epochs. A major challenge lies in the absence of a metric to quantify how well the pretrained generative ability is preserved. Current finetuning techniques implicitly assume that a smaller Euclidean distance between the finetuned and pretrained models signifies better preservation of pretrained capabilities. Consequently, finetuning is usually conducted with a very low learning rate. While this assumption can sometimes be valid, we contend that relying solely on Euclidean distance to the pretrained model fails to fully capture the extent of semantic preservation. Therefore, employing a more structural metric to characterize the differences between the finetuned and pretrained models could significantly enhance the preservation of pretrained performance and improve finetuning stability.

Motivated by empirical findings that hyperspherical similarity effectively captures semantic information~\cite{liu2017deep,liu2018decoupled,chen2020angular}, we utilize hyperspherical energy~\cite{liu2018learning} to describe the pairwise relational structure among neurons. Hyperspherical energy is defined as the sum of hyperspherical similarities (e.g., cosine similarity) between every pair of neurons within the same layer, reflecting the degree of neuron uniformity on the unit hypersphere~\cite{liu2021learning}. We propose that an optimal finetuned model should exhibit minimal deviation in hyperspherical energy relative to the pretrained model. A straightforward approach would be to incorporate a regularizer that maintains constant hyperspherical energy during finetuning; however, this does not guarantee effective minimization of the energy difference. To address this, we leverage an invariance property of hyperspherical energy—namely, that pairwise hyperspherical similarity is exactly preserved when the same orthogonal transformation is applied to all neurons. Building on this property, we introduce Orthogonal Finetuning (OFT), which adapts large text-to-image diffusion models to downstream tasks without altering their hyperspherical energy. The key concept is to learn a shared orthogonal transformation per layer that preserves the pairwise angles among neurons. OFT can also be interpreted as modifying the canonical coordinate system within each layer. Considering that smaller Euclidean distances between the finetuned and pretrained models suggest better retention of pretrained performance, we further propose a variant called Constrained Orthogonal Finetuning (COFT), which restricts the finetuned model to lie on a hypersphere of fixed radius centered around the pretrained neurons.

The rationale behind why orthogonal transformation is effective for finetuning neurons partially derives from the 2D Fourier transform, which decomposes an image into magnitude and phase spectra. The phase spectrum, representing the angular relationship between input and basis, retains most of the semantic content. For instance, an image's phase spectrum combined with a random magnitude spectrum can still reproduce the original image without losing its semantics. This observation implies that altering neuron directions is crucial for semantically modifying the generated image—an objective shared by both subject-driven and controllable generation. However, allowing unrestricted changes in neuron directions inevitably harms the pretrained generative performance. To limit this freedom, we propose maintaining the angle between every pair of neurons, based on the premise that inter-neuronal angles are essential for encoding neural network knowledge. Given this insight, it is natural to learn a layer-shared orthogonal transformation for neurons within each layer to keep the hyperspherical energy constant.

Our approach is also inspired by orthogonal over-parameterized training~\cite{liu2021orthogonal}, which trains classification networks from scratch by orthogonally transforming a randomly initialized neural network. This method relies on the fact that a randomly initialized network exhibits, on average, low hyperspherical energy, and the objective of \cite{liu2021orthogonal} is to maintain low hyperspherical energy throughout training (since low energy correlates with better classification generalization~\cite{liu2018learning,lin2020regularizing}). The study \cite{liu2021orthogonal} demonstrates that orthogonal transformation is sufficiently expressive to train generalizable classification networks. In contrast, our focus lies in finetuning text-to-image diffusion models to achieve enhanced controllability and improved downstream generative quality. We highlight two main distinctions between OFT and \cite{liu2021orthogonal}. First, while \cite{liu2021orthogonal} aims to minimize hyperspherical energy, OFT strives to preserve the hyperspherical energy of the pretrained model, ensuring that the intrinsic pretrained structure remains intact during finetuning. Minimizing hyperspherical energy in diffusion model finetuning could disrupt original semantic structures. Second, OFT minimizes deviation from the pretrained model, leading to a constrained formulation, whereas \cite{liu2021orthogonal} imposes no such restrictions. Striking an effective balance between flexibility and stability is crucial for finetuning, and we contend that our OFT framework successfully accomplishes this. Our key contributions are summarized as follows:

\begin{itemize}[leftmargin=*,nosep]

    \item We introduce a new finetuning technique—Orthogonal Finetuning (OFT)—to enhance controllability in text-to-image diffusion models. To increase stability, we also propose a constrained version that restricts the angular deviation from the pretrained model.
    \item Unlike existing finetuning approaches, OFT fine-tunes the model while provably maintaining the hyperspherical energy, which we empirically identify as a crucial indicator for preserving the generative semantics of the pretrained model.
    \item We utilize OFT in two applications: subject-driven generation and controllable generation. Through extensive empirical evaluation, we show considerable improvements over previous methods in terms of generation quality, convergence speed, and finetuning robustness. Additionally, OFT demonstrates superior sample efficiency, achieving convergence with substantially fewer training images and epochs.
    \item For controllable generation, we propose a novel control consistency metric to assess controllability. The main idea involves estimating the control signal from the generated image and comparing it to the original control signal. This metric further confirms OFT's strong controllability.
\end{itemize}

\section{Related Work}

\textbf{Text-to-image diffusion models}. Substantial advancements~\cite{saharia2022photorealistic,ramesh2022hierarchical,rombach2022high,gu2022vector,nichol2022glide} have been achieved in text-to-image generation, largely driven by the rapid progress in diffusion-based generative models~\cite{ho2020denoising,song2021score,song2021denoising,dhariwal2021diffusion} and vision-language representation learning~\cite{su2020vlbert,lu2019vilbert,radford2021learning,wang2021simvlm,li2022blip,li2023blip,alayrac2022flamingo,schuhmann2022laion}. GLIDE~\cite{nichol2022glide} and Imagen~\cite{saharia2022photorealistic} are trained directly in pixel space. GLIDE jointly trains the text encoder with a diffusion prior using paired text-image datasets, whereas Imagen employs a frozen pretrained text encoder. Stable Diffusion~\cite{rombach2022high} and DALL-E2~\cite{ramesh2022hierarchical} operate in latent space. Stable Diffusion uses VQ-GAN~\cite{esser2021taming} to learn a latent visual codebook, while DALL-E2 leverages CLIP~\cite{radford2021learning} to create a joint latent embedding space for images and text. Beyond diffusion models, generative adversarial networks~\cite{reed2016generative,zhang2017stackgan,xu2018attngan,li2019controllable} and autoregressive models~\cite{ramesh2021zero,ding2021cogview,wu2022nuwa,yuscaling} have also been explored for text-to-image generation. OFT is fundamentally a model-agnostic finetuning strategy applicable to any text-to-image diffusion model.

\textbf{Subject-driven generation}. To avoid altering the subject, \cite{avrahami2022blended,nichol2022glide} use user-provided masks as additional conditions. Inversion methods~\cite{choi2021ilvr,dhariwal2021diffusion,ramesh2022hierarchical,gal2022image} modify context without changing the subject. \cite{hertz2022prompt} supports both local and global edits without requiring input masks. However, these approaches often struggle to retain identity-related subject details. Pivotal Tuning~\cite{roich2022pivotal} finetunes the generator around an initial inverted latent code with added regularization to preserve identity. Similarly, \cite{nitzan2022mystyle} learns a personalized generative face prior from multiple images of a person. \cite{casanova2021instance} generates diverse variations of an instance but may lose instance-specific details. DreamBooth~\cite{ruiz2023dreambooth} finetunes a text-to-image diffusion model with a customized token and limited subject images, using a reconstruction loss combined with a class-specific prior preservation loss. OFT builds on the DreamBooth framework but replaces naive low learning rate finetuning with orthogonal transformation-based finetuning.

\textbf{Controllable generation}. Image-to-image translation can be regarded as a type of controllable generation, with many prior approaches mainly employing conditional generative adversarial networks~\cite{isola2017image,zhu2017toward,wang2018high,park2019semantic,choi2018stargan}. Diffusion models have also been applied to this task~\cite{saharia2022palette,voynov2022sketch,wang2022pretraining}. Recently, ControlNet~\cite{zhang2023adding} introduced a method to steer a pretrained diffusion model by finetuning it to incorporate additional control signals, achieving remarkable performance in controllable generation. Another contemporaneous and related work, T2I-Adapter~\cite{mou2023t2i}, similarly finetunes pretrained diffusion models to enhance the controllability of the generated images. Adopting the same task setting as in \cite{zhang2023adding,mou2023t2i}, we utilize OFT to finetune pretrained diffusion models, resulting in consistently improved controllability with fewer training samples and reduced finetuning parameters. Notably, OFT does not add any extra computational cost during inference.

\textbf{Model finetuning}. The practice of finetuning large pretrained models on downstream tasks has gained significant traction recently~\cite{devlin2018bert,brown2020language,he2022masked}. Within the context of finetuning, adaptation techniques (\emph{e.g.}, \cite{hulora2022,houlsby2019parameter,pfeiffer2020adapterfusion}) have been extensively explored in natural language processing. LoRA~\cite{hulora2022} is particularly relevant to OFT, as it assumes a low-rank form for additive weight updates during finetuning. In contrast, OFT applies a layer-shared orthogonal transformation that multiplicatively updates neuron weights and provably maintains pairwise angles among neurons within the same layer, resulting in enhanced stability.

\section{Orthogonal Finetuning}

\subsection{Why Does Orthogonal Transformation Make Sense?}

We begin by examining why orthogonal transformation is advantageous for finetuning text-to-image diffusion models. This question is broken down into two parts: (1) why it is beneficial to finetune neuron angles (\emph{i.e.}, directions), and (2) why orthogonal transformations are employed to adjust these angles.

Regarding the first question, we take motivation from the empirical findings in \cite{liu2018decoupled,chen2020angular} which indicate that the angular difference in features effectively captures the semantic gap. SphereNet~\cite{liu2017deep} demonstrates that training a neural network with all neurons normalized to lie on a unit hypersphere maintains comparable capacity and even improves generalizability, suggesting that the direction of neurons can fully encapsulate the most critical information from the data. To further highlight the significance of neuron angles, we conduct a simple experiment shown in Figure~\ref{fig:toy_ae} where a standard convolutional autoencoder is trained on a set of flower images. During training, the feature map is produced using the conventional inner product ($z$ represents the output element of the convolution kernel $\bm{w}$ and $\bm{x}$ is the input within the sliding window). At test time, we compare three methods for generating the feature map: (a) the inner product as used during training, (b) the magnitude information, and (c) the angular information. The results in Figure~\ref{fig:toy_ae} illustrate that the angular information of neurons can almost perfectly reconstruct the input images, whereas the magnitude of neurons lacks useful information. It is important to note that cosine activation (c) is not applied during training; training is conducted solely based on the inner product. These findings suggest that neuron angles (directions) primarily encode the semantic content of input images. Thus, adjusting neuron directions is likely more effective for modifying the semantic information in images.

Concerning the second question, the most straightforward method to fine-tune neuron directions is by simultaneously rotating or reflecting all neurons within the same layer, which naturally leads to orthogonal transformations. Although more flexible approaches might involve different angular rotations for individual neurons, we find orthogonal transformations strike an optimal balance between flexibility and structure. Additionally, \cite{liu2021orthogonal} demonstrates that orthogonal transformations possess sufficient expressive power for neural network learning. To support this, we conduct an experiment illustrating the beneficial regularization effect of the orthogonality constraint. Using the ControlNet~\cite{zhang2023adding} setup, we perform a controllable generation task, with results shown in Figure~\ref{fig:ortho}. Our standard OFT method demonstrates stable performance and achieves precise control after training (epoch 20). By contrast, OFT without the orthogonality constraint fails to generate realistic images and does not achieve any control effect. This experiment underscores the critical role of the orthogonality constraint in OFT.

\subsection{General Framework}

The traditional finetuning approach generally employs gradient descent with a small learning rate to update either the entire model or specific layers. This small learning rate implicitly restricts the updated model to remain close to the pretrained model, so standard finetuning effectively aims to train the model while implicitly minimizing $\thickmuskip=2mu \medmuskip=2mu \| \bm{M} - \bm{M}^0\|$, where $\bm{M}$ represents the finetuned model weights and $\bm{M}^0$ denotes the pretrained weights. Although this implicit constraint is intuitively reasonable, it may still allow excessive flexibility when finetuning large models. To overcome this limitation, LoRA imposes an additional low-rank constraint on the weight update, specifically $\thickmuskip=2mu \medmuskip=2mu \text{rank}(\bm{M} - \bm{M}^0)=r'$, where $r'$ is chosen to be a small value. In contrast, OFT introduces a constraint based on pairwise neuron similarity: $\thickmuskip=2mu \medmuskip=2mu \| \text{HE}(\bm{M})-\text{HE}(\bm{M}^0) \|=0$, where $\text{HE}(\cdot)$ denotes the hyperspherical energy of a weight matrix. To illustrate, consider a fully connected layer $\thickmuskip=2mu \medmuskip=2mu \bm{W}=\{\bm{w}_1,\cdots,\bm{w}_n\}\in\mathbb{R}^{d\times n}$ with neurons $\bm{w}_i \in \mathbb{R}^d$ for $i$-th neuron ($\bm{W}^0$ being the pretrained weights). The output vector $\thickmuskip=2mu \medmuskip=2mu \bm{z}\in\mathbb{R}^n$ of this layer is computed as $\thickmuskip=2mu \medmuskip=2mu \bm{z}=\bm{W}^\top\bm{x}$ where $\thickmuskip=2mu \medmuskip=2mu \bm{x}\in\mathbb{R}^d$ is the input vector. OFT can be understood as minimizing the difference in hyperspherical energy between the finetuned and pretrained models:

\begin{equation}\label{eq:oft_principle}
\footnotesize
\min_{\bm{W}} \left\| \text{HE}(\bm{W})-\text{HE}(\bm{W}^0)\right\|~~\Leftrightarrow~~\min_{\bm{W}} \bigg{\|} \sum_{i\neq j} \|\hat{\bm{w}}_i-\hat{\bm{w}}_j\|^{-1}-\sum_{i\neq j} \|\hat{\bm{w}}^0_i-\hat{\bm{w}}^0_j\|^{-1}\bigg{\|}
\end{equation}

Here, $\thickmuskip=2mu \medmuskip=2mu \hat{\bm{w}}_i:=\bm{w}_i/\|\bm{w}_i\|$ denotes the normalized $i$-th neuron, and the hyperspherical energy for a layer $\bm{W}$ is defined as $\thickmuskip=2mu \medmuskip=2mu \text{HE}(\bm{W}):= \sum_{i\neq j} \|\hat{\bm{w}}_i-\hat{\bm{w}}_j\|^{-1}$. It is straightforward to see that the minimum value of zero can be achieved in Eq.~\eqref{eq:oft_principle}. This minimum is attained when $\bm{W}$ and $\bm{W}^0$ differ only by a rotation or reflection, i.e., $\thickmuskip=2mu \medmuskip=2mu \bm{W}=\bm{R}\bm{W}^0$, where $\thickmuskip=2mu \medmuskip=2mu \bm{R}\in\mathbb{R}^{d\times d}$ is an orthogonal matrix (with determinant $1$ indicating rotation and $-1$ indicating reflection). This concept embodies the core of OFT: the finetuning process only involves learning layer-shared orthogonal transformations to map neurons in each layer. Formally, OFT seeks to

\begin{equation}\label{eq:oft_general}
    \footnotesize
    \bm{z}=\bm{W}^\top\bm{x}=(\bm{R}\cdot\bm{W}^0)^\top\bm{x},~~~\text{s.t.}~\bm{R}^\top\bm{R}=\bm{R}\bm{R}^\top=\bm{I}
\end{equation}

where $\bm{W}$ represents the OFT-finetuned weight matrix and $\bm{I}$ denotes the identity matrix. The concept of OFT is depicted in Figure~\ref{fig:block}. Analogous to the zero initialization used in LoRA, it is necessary to ensure that OFT fine-tunes the pretrained model precisely starting from $\bm{W}^0$. To accomplish this, we initialize the orthogonal matrix $\bm{R}$ as the identity matrix, so that the fine-tuned model begins with the pretrained weights. To maintain the orthogonality of the matrix $\bm{R}$, differential orthogonalization methods described in \cite{liu2021orthogonal,lezcano2019cheap} can be utilized. We will elaborate on a computationally efficient approach to preserving orthogonality.

\subsection{Efficient Orthogonal Parameterization}\label{sect:eff_ortho}

Although standard orthogonalization techniques like the Gram-Schmidt process are differentiable, they are often computationally costly in practical scenarios~\cite{liu2021orthogonal}. To improve efficiency, we employ Cayley parameterization to generate the orthogonal matrix. Specifically, the orthogonal matrix is constructed as $\thickmuskip=2mu \medmuskip=2mu \bm{R}=(\bm{I}+\bm{Q})(\bm{I}-\bm{Q})^{-1}$, where $\bm{Q}$ is a skew-symmetric matrix satisfying $\thickmuskip=2mu \medmuskip=2mu \bm{Q}=-\bm{Q}^\top$. This approach offers computational benefits but has a minor limitation: the Cayley parameterization can only yield orthogonal matrices with determinant $1$, belonging to the special orthogonal group. Fortunately, we observe that this restriction does not impact practical performance. Even when using the Cayley transform to parameterize $\bm{R}$, it can still become parameter-inefficient for large dimension $d$. To mitigate this, we propose representing $\bm{R}$ as a block-diagonal matrix composed of $r$ blocks, resulting in the following structure:

\begin{equation}
    \footnotesize
    \bm{R}=\text{diag}(\bm{R}_1,\bm{R}_2,\cdots,\bm{R}_r)=\begin{bmatrix}
    \bm{R}_1\in \text{O}\left(\frac{d}{r}\right) &  &\\
    &\ddots & \\
    & & \bm{R}_r\in \text{O}\left(\frac{d}{r}\right)
    \end{bmatrix}\in \text{O}(d)
\end{equation}

where $\text{O}(d)$ represents the orthogonal group in dimension $d$, and $\thickmuskip=2mu \medmuskip=2mu \bm{R}\in\mathbb{R}^{d\times d}$ along with $\thickmuskip=2mu \medmuskip=2mu \bm{R}_i\in\mathbb{R}^{d/r\times d/r},\forall i$ are orthogonal matrices. When $\thickmuskip=2mu \medmuskip=2mu r=1$, the block-diagonal orthogonal matrix reduces to a conventional unconstrained orthogonal matrix. For an orthogonal matrix of size $\thickmuskip=2mu \medmuskip=2mu d\times d$, the parameter count is $\thickmuskip=2mu \medmuskip=2mu d(d-1)/2$, which leads to a computational complexity of $\mathcal{O}(d^2)$. For an $r$-block diagonal orthogonal matrix, the parameter count is $\thickmuskip=2mu \medmuskip=2mu d(d/r-1)/2$, resulting in complexity $\thickmuskip=2mu \medmuskip=2mu \mathcal{O}(d^2/r)$. One may further decrease the number of parameters by sharing the block matrices, \emph{i.e.}, setting $\thickmuskip=2mu \medmuskip=2mu \bm{R}_i=\bm{R}_j, \forall i \neq j$. This reduces parameter complexity to $\mathcal{O}(d^2/r^2)$. Despite these parameter efficiency improvements, the resulting matrix $\bm{R}$ retains its orthogonality, thus there is no loss in preserving hyperspherical energy.

Next, we compare the parameter efficiency of OFT with LoRA. For LoRA with a low-rank parameter $r'$, the number of trainable parameters is $\thickmuskip=2mu \medmuskip=2mu r'(d+n)$. Assuming both $r$ and $r'$ depend on the neuron dimension $d$ (for example, $\thickmuskip=2mu \medmuskip=2mu r = r' = \alpha d$ where $\thickmuskip=2mu \medmuskip=2mu 0 < \alpha \leq 1$ is a constant), the parameter complexity of LoRA is $\mathcal{O}(d^2 + dn)$, whereas for OFT it is $\mathcal{O}(d)$. To illustrate the complexity difference concretely, consider a weight matrix of size $\thickmuskip=2mu \medmuskip=2mu 128 \times 128$: LoRA has $2,048$ trainable parameters with $\thickmuskip=2mu \medmuskip=2mu r' = 8$, whereas OFT has $960$ trainable parameters with $\thickmuskip=2mu \medmuskip=2mu r = 8$ (without applying block sharing).

\subsection{Constrained Orthogonal Finetuning}

The original OFT can be further restricted by confining the finetuned model within a small vicinity of the pretrained model. In particular, COFT employs the forward pass:

\begin{equation}\label{eq:coft_general}
    \footnotesize
    \bm{z} = \bm{W}^\top \bm{x} = (\bm{R} \cdot \bm{W}^0)^\top \bm{x}, \quad \text{s.t.} \quad \bm{R}^\top \bm{R} = \bm{R} \bm{R}^\top = \bm{I}, \quad \|\bm{R} - \bm{I}\| \leq \epsilon
\end{equation}

which imposes both an orthogonality constraint and an $\epsilon$-bounded deviation from the identity matrix. The orthogonality constraint can be enforced using the Cayley parameterization introduced in Section~\ref{sect:eff_ortho}. However, incorporating the $\epsilon$-deviation constraint into the Cayley-parameterized orthogonal matrix is not straightforward. To better understand the Cayley transform, we apply the Neumann series to approximate $\thickmuskip=2mu \medmuskip=2mu \bm{R} = (\bm{I} + \bm{Q})(\bm{I} - \bm{Q})^{-1}$ as $\th

series converges in the operator norm). Consequently, the constraint $\thickmuskip=2mu \medmuskip=2mu\|\bm{R}-\bm{I}\|\leq\epsilon$ can be incorporated within the Cayley transform, leading to an equivalent constraint $\thickmuskip=2mu \medmuskip=2mu \|\bm{Q}-\bm{0}\|\leq\epsilon'$, where $\bm{0}$ is the zero matrix and $\epsilon'$ is a distinct error hyperparameter (different from $\epsilon$). This new restriction on the matrix $\bm{Q}$ can be straightforwardly enforced via projected gradient descent. To initialize the orthogonal matrix $\bm{R}$ as the identity, we start with $\bm{Q}$ as the zero matrix. COFT can be interpreted as combining two explicit constraints: a minimal hyperspherical energy difference and a bounded deviation from the pretrained model. While the latter constraint is often implicitly assumed in existing finetuning approaches, COFT explicitly incorporates it. Although OFT demonstrates strong performance, we find that COFT provides enhanced finetuning stability compared to OFT, attributed to the explicit deviation constraint. Figure~\ref{fig:eps_coft} illustrates how the parameter $\epsilon$ influences COFT performance. It is apparent that $\epsilon$ governs the finetuning flexibility: larger $\epsilon$ makes the COFT-finetuned model resemble the OFT-finetuned one more closely, whereas smaller $\epsilon$ causes the COFT-finetuned model to behave increasingly like the original pretrained text-to-image diffusion model.

\subsection{Re-scaled Orthogonal Finetuning}

We introduce a straightforward extension to the original OFT by additionally learning a scaling coefficient for each neuron's magnitude. This idea stems from the observation that rescaling neurons does not affect hyperspherical energy because magnitude normalization removes its influence. Specifically, the forward pass is given by: $\thickmuskip=2mu \medmuskip=2mu\bm{z}=(\bm{R}\bm{W}^0\bm{D})^\top\bm{x}$\footnote{\textbf{Errata}: The NeurIPS camera-ready version incorrectly states the re-scaled OFT forward pass as $\thickmuskip=2mu \medmuskip=2mu\bm{z}=(\bm{D}\bm{R}\bm{W}^0)^\top\bm{x}$. The implementation used was correct, so the results reported in Appendix~\ref{app:rescaled} remain valid.}, where $\thickmuskip=2mu \medmuskip=2mu \bm{D}=\text{diag}(s_1,\cdots,s_n)\in\mathbb{R}^{n\times n}$ is a learnable diagonal matrix with strictly positive diagonal elements $s_1,\cdots,s_n$. Unlike the original OFT forward pass in Eq.~\eqref{eq:oft_general}, where only $\bm{R}$ is learnable, here both the diagonal matrix $\bm{D}$ and the orthogonal matrix $\bm{R}$ are trainable. This re-scaled OFT enhances the flexibility of OFT with negligible additional parameter cost. While experiments are conducted with the original OFT to isolate the effectiveness of orthogonal transformation alone, we observe that the re-scaled OFT typically performs better (see Appendix~\ref{app:rescaled}).

\section{Intriguing Insights and Discussions}

\textbf{OFT is architecture-agnostic}. In principle, OFT can be applied to any neural network architecture. For Transformers, LoRA is commonly applied to attention weights~\cite{hulora2022}. To enable fair comparison with LoRA, we restrict OFT finetuning to the attention weights in our experiments. Beyond fully connected layers, OFT is also well-suited for finetuning convolutional layers because the block-diagonal structure of $\bm{R}$ has meaningful interpretations in this context (unlike LoRA). When the number of blocks matches the number of input channels, each block transforms a single neuron channel independently, akin to learning depth-wise convolution kernels~\cite{chollet2017xception}. If all blocks in $\bm{R}$ are shared, OFT applies a channel-shared orthogonal transformation to the neurons. A preliminary study on finetuning convolution layers with OFT is provided in Appendix~\ref{app:convolution}.

\textbf{Relation to LoRA}. LoRA adds a low-rank matrix to avoid significant alterations to the pretrained weight matrix's information. Conversely, OFT enforces an orthogonal (full-rank) transformation on the pretrained weight matrix, ensuring that the transformation preserves the pretrained knowledge. The forward pass of OFT can be expressed as $\thickmuskip=2mu \medmuskip=2mu \bm{z}=(\bm{R}\bm{W}^0)^\top\bm{x}=(\bm{W}^0+(\bm{R}-\bm{I})\bm{W}^0)^\top\bm{x}$, where $\thickmuskip=2mu \medmuskip=2mu (\bm{R}-\bm{I})\bm{W}^0$ is analogous to the low-rank weight update used in LoRA. Given that $\bm{W}^0$ is generally full-rank, OFT effectively performs a low-rank update when $\thickmuskip=2mu \medmuskip=2mu \bm{R}-\bm{I}$ is low-rank. Similar to how LoRA uses a rank parameter $r'$ to limit trainable parameters, OFT employs a diagonal block parameter $r$ for parameter efficiency. Interestingly, LoRA and OFT embody two fundamentally different approaches to parameter efficiency: LoRA leverages low-rank structure, whereas OFT capitalizes on sparsity patterns, specifically block-diagonal orthogonality.

\textbf{Reasons for OFT's faster convergence}. First, as illustrated in Figure~\ref{fig:toy_ae}, the most impactful updates for altering semantic representations are those that modify neuron directions, which aligns directly with OFT’s design. Additionally, OFT can be interpreted as fine-tuning neurons constrained on a smooth hypersphere manifold, resulting in an improved optimization landscape. This has also been empirically demonstrated in \cite{liu2021orthogonal}.

\textbf{Why we do not minimize hyperspherical energy}. Unlike \cite{liu2021orthogonal}, our approach does not aim to minimize hyperspherical energy. In classification tasks, having neurons without redundancy is preferable. Minimizing hyperspherical energy would space neurons evenly across the hypersphere, but this objective is not suitable for fine-tuning because it risks disrupting the information acquired during pretraining.

\textbf{Balancing flexibility and regularity in finetuning}. We identify a fundamental trade-off between flexibility and regularity. Standard finetuning offers the greatest flexibility but often results in poor stability and a higher risk of model collapse. Despite its simplicity, OFT achieves an effective compromise between flexibility and regularity by maintaining the pairwise neuron angles. The block-diagonal parameterization can also be interpreted as a stronger form of regularization applied to the orthogonal matrix.

\textbf{No extra inference cost}. Unlike ControlNet, our OFT framework does not add any extra computational burden during inference for the finetuned model. At inference time, we simply multiply the learned orthogonal matrix $\bm{R}$ with the pretrained weight matrix $\bm{W}^0$ to obtain an equivalent weight matrix $\thickmuskip=2mu \medmuskip=2mu \bm{W}=\bm{R}\bm{W}^0$. Consequently, the inference speed remains identical to that of the pretrained model.

\section{Experiments and Results}

\textbf{Overall setup}. In our experiments, we use Stable Diffusion v1.5~\cite{rombach2022high} as the pretrained text-to-image model. To ensure fairness, we randomly select generated images from each method. For subject-driven generation, we broadly follow the DreamBooth~\cite{ruiz2023dreambooth} approach. For controllable generation, we mostly adhere to ControlNet~\cite{zhang2023adding} and T2I-Adapter~\cite{mou2023t2i}. To maintain a fair comparison with LoRA, we restrict the application of OFT or COFT to the same layers where LoRA is applied. Additional results and specifics are provided in Appendix~\ref{app:settings}.

\subsection{Subject-driven Generation}

\textbf{Experimental setup}. We compare against DreamBooth~\cite{ruiz2023dreambooth} and LoRA~\cite{hulora2022} as baseline methods. All methods utilize the same loss function as DreamBooth. For DreamBooth and LoRA, we generally follow the original papers and adopt the optimal hyperparameter configurations. Further results can be found in Appendix~\ref{app:settings},\ref{app:coft_oft},\ref{app:more_qualitative},\ref{app:failure}.

\textbf{Stability and convergence of finetuning}. We initially assess the finetuning stability and speed of convergence for DreamBooth, LoRA, OFT, and COFT. The results are displayed in Figure~\ref{fig:teaser} and Figure~\ref{fig:db_convergence}. It is evident that both COFT and OFT can finetune the diffusion model in a stable manner. After 400 iterations, DreamBooth and OFT variants both demonstrate strong control, whereas LoRA fails to maintain the subject's identity. By 2000 iterations, DreamBooth begins producing collapsed images, and LoRA fails to generate the yellow shirt correctly, instead producing yellow fur. In contrast, OFT and COFT continue to maintain stable and consistent control over the output images. These findings confirm the rapid yet stable convergence of our OFT framework in subject-driven generation tasks. We emphasize that being insensitive to the number of finetuning iterations is crucial, as it reduces the difficulty of tuning iteration counts for various subjects. For both OFT and COFT, it is possible to select a relatively high iteration number without extensive tuning. In particular, for COFT with an appropriate $\epsilon$, setting both the learning rate and iteration count becomes straightforward.

\textbf{Quantitative evaluation}. In line with \cite{ruiz2023dreambooth}, we perform a quantitative comparison to assess subject fidelity (using DINO~\cite{caron2021emerging}, CLIP-I~\cite{radford2021learning}), text prompt fidelity (CLIP-T~\cite{radford2021learning}), and sample diversity (LPIPS~\cite{zhang2018unreasonable}). CLIP-I measures the average pairwise cosine similarity of CLIP embeddings between generated and real images. DINO is analogous to CLIP-I but utilizes ViT S/16 DINO embeddings instead. CLIP-T calculates the average cosine similarity between text prompts and generated images' CLIP embeddings. We further compute the average LPIPS cosine similarity among generated images of the same subject given identical text prompts. Table~\ref{table:dreambooth_final} indicates that COFT and OFT substantially outperform DreamBooth and LoRA on the DINO and CLIP-I metrics, while delivering slightly improved or comparable outcomes in prompt fidelity and diversity. To reduce randomness, each method finetunes the same pretrained model 30 times with different random seeds. The results demonstrate that our OFT framework not only offers superior convergence and stability but also consistently achieves better final results.

\textbf{Qualitative comparison}. To gain a clearer and more intuitive insight into the advantages of OFT, we present a selection of randomly chosen examples for subject-driven generation in Figure~\ref{fig:dreambooth_final}. For an unbiased comparison, all examples are produced from the identical finetuned model via each approach, ensuring that no individual text prompt is optimized specifically for the final outcomes. For each approach, we pick the model that attains the highest validation CLIP scores. The results displayed in Figure~\ref{fig:dreambooth_final} reveal that both OFT and COFT provide excellent preservation of semantic subject identity, whereas LoRA frequently struggles to maintain the subject's identity (for instance, LoRA completely loses the subject identity in the bowl example). Simultaneously, OFT and COFT demonstrate much finer control through text prompts, while DreamBooth, despite preserving subject identity, often fails to produce images aligned with the text prompts (e.g., the first row of the bowl example). This qualitative comparison highlights that our OFT framework achieves superior control and subject retention simultaneously. Additionally, OFT's performance is robust to the number of iterations, maintaining effectiveness even with many iterations, unlike DreamBooth and LoRA. Further qualitative examples are provided in Appendix~\ref{app:more_qualitative}. We also perform a human evaluation in Appendix~\ref{app:human}, which further confirms the advantages of OFT.

\subsection{Controllable Generation}

\textbf{Settings}. We use ControlNet~\cite{zhang2023adding}, T2I-Adapter~\cite{mou2023t2i}, and LoRA~\cite{hulora2022} as baseline methods. We focus on three challenging controllable generation tasks in the main paper: Canny edge to image (C2I) on the COCO dataset~\cite{lin2014microsoft}, segmentation map to image (S2I) on the ADE20K dataset~\cite{zhou2017scene}, and landmark to face (L2F) on the CelebA-HQ dataset~\cite{karras2018progressive,xia2021tedigan}. All methods are applied to finetune Stable Diffusion (SD) v1.5 on these datasets for 20 epochs. Additional results can be found in Appendices~\ref{app:more_qualitative}, \ref{app:more_control_task}, and \ref{app:failure}.

\textbf{Convergence}. We assess the convergence rates of ControlNet, T2I-Adapter, LoRA, and COFT on the L2F task through both quantitative and qualitative analyses. For the quantitative metric, we calculate the average $\ell_2$ distance between the control face landmarks and the predicted face landmarks. In Figure~\ref{fig:face_convergence}, we display the face landmark error for models finetuned over various numbers of epochs. It is evident that both COFT and OFT converge much more rapidly. LoRA requires 20 epochs to reach the performance level that our OFT framework achieves by the 8th epoch. It is noteworthy that OFT and COFT employ a comparable number of trainable parameters to LoRA (considerably fewer than ControlNet), yet they converge far more efficiently than existing approaches. Additionally, the swift convergence of OFT is further supported by the results shown in Figure~\ref{fig:teaser}. The example on the right in Figure~\ref{fig:teaser} illustrates that OFT is substantially more data-efficient than ControlNet and LoRA, as OFT achieves good convergence using only 5\% of the ADE20K dataset. For qualitative comparisons, we concentrate on OFT, COFT, and ControlNet, since ControlNet attains the closest landmark error to our methods. The outcomes in Figure~\ref{fig:face_convergence_qual} demonstrate that both OFT and COFT exhibit stable convergence, with the generated face pose progressively aligning with the control landmarks. In contrast, ControlNet’s controllability abruptly appears after the 8th epoch, reflecting the sudden convergence pattern reported in \cite{zhang2023adding}. This stability in convergence makes our OFT framework considerably easier to apply in practice, as OFT’s training dynamics are much smoother and more predictable. Consequently, it becomes simpler to tune OFT’s hyperparameters effectively.

\textbf{Quantitative comparison}. We propose a control consistency metric to assess the effectiveness of controllable generation. The core concept is to extract the control signal from the generated image and then compare it against the original input control signal. For the C2I task, we calculate IoU and F1 score. For the S2I task, we measure mean IoU, mean accuracy, and overall accuracy. For the L2F task, we determine the mean $\ell_2$ distance between the control landmarks and the predicted landmarks. Additional details about the consistency metrics are provided in Appendix~\ref{app:settings}. For all methods compared, we utilize the best achievable hyperparameter configurations. The results displayed in Table~\ref{table:control_gen} demonstrate that both OFT and COFT produce significantly stronger and more precise control compared to other approaches. We notice that adapter-based methods (e.g., T2I-Adapter and ControlNet) tend to converge more slowly and yield inferior final outcomes. Relative to ControlNet, LoRA achieves better performance on the S2I task but worse results on the C2I and L2F tasks. Overall, we observe that LoRA’s performance limit is relatively low, even after careful hyperparameter tuning. In contrast, the performance of our OFT framework has not reached saturation, as it empirically improves with an increasing number of trainable parameters. We highlight that our quantitative evaluation of controllable generation represents a novel contribution, as it enables precise assessment of the control capabilities of finetuned models (subject to the accuracy of the utilized off-the-shelf segmentation/detection model).

\textbf{Qualitative comparison}. We also conduct a qualitative comparison between OFT, COFT, and leading contemporary methods such as ControlNet, T2I-Adapter, and LoRA. The randomly generated images in Figure~\ref{fig:control_final} illustrate that OFT and COFT not only produce high-fidelity and realistic images but also achieve precise control. In the S2I task, it is evident that LoRA completely fails to generate images that adhere to the input segmentation map, whereas ControlNet, OFT, and COFT successfully control the generated images. Compared to ControlNet, both OFT and COFT generate higher-fidelity images with more vivid details and more plausible geometric structures while using significantly fewer model parameters. For the C2I task, OFT and COFT are capable of hallucinating semantically similar images based on rough Canny edges, whereas T2I-Adapter and LoRA show markedly inferior performance. In the L2F task, our approach achieves the most accurate pose control for generated faces, even under challenging facial poses. Across all three control tasks, we demonstrate that OFT and COFT produce qualitatively superior images compared to state-of-the-art baselines, highlighting the effectiveness of our OFT framework for controllable image generation. To provide a more extensive qualitative comparison, additional examples for all three control tasks are included in Appendix~\ref{app:control_exp}. Furthermore, we show that OFT performs well on additional control tasks (such as dense pose to human body, sketch to image, and depth to image) in Appendix~\ref{app:more_control_task}.

\section{Concluding Remarks and Open Problems}

Inspired by the insight that angular relationships between neurons play a key role in defining visual semantics, we introduce a straightforward yet powerful finetuning technique—orthogonal finetuning (OFT)—for managing text-to-image diffusion models. Our focus lies on two specific applications: subject-driven generation and controllable generation. Compared with prior approaches, OFT offers enhanced controllability and finetuning stability while requiring fewer finetuning parameters. Crucially, OFT imposes no additional computational cost during inference, making it a highly efficient and deployable solution.

OFT also raises several intriguing open questions. Firstly, the orthogonality in OFT is ensured through Cayley parametrization, which necessitates a matrix inversion step. This operation somewhat constrains OFT’s scalability. Although we partially mitigate this limitation by adopting a block diagonal parametrization, accelerating the matrix inversion in a differentiable manner remains an open challenge. Secondly, OFT exhibits promising compositional properties, as orthogonal matrices derived from multiple OFT finetuning tasks can be multiplied to produce another orthogonal matrix. Investigating whether this set of orthogonal matrices retains the knowledge from all downstream tasks presents an engaging avenue for future research. Lastly, the parameter efficiency of OFT heavily relies on the block diagonal structure, which inherently introduces biases and restricts flexibility. Developing approaches to enhance parameter efficiency in a more effective and less biased manner stands out as a significant unresolved problem.

\end{document}